\chapter*{Conclusions et perspectives}
\addcontentsline{toc}{chapter}{Conclusions et perspectives}  

Le dévelopement actuel de la robotique tend à indiquer que les interactions physiques des robots avec des humains et leurs environnements vont se multiplier. En s'inspirant des propriétés biomécaniques, et aussi neuromusculaires observées chez l'Homme, de nombreux chercheurs continuent de proposer des structures et stratégies de contrôles innovantes pour améliorer les performances et la robustesse des machines, notamment aux abords et contacts d'humains. 

Une partie de la recherche se consacre ainsi à l'étude sociale des interactions entre humains, et aussi machines, afin de rendre leur présence et déplacements plus naturels et leur comportement plus acceptable et prévisible. Tandis qu'une autre partie de la recherche s'intéresse principalement aux performances et à la robustesse des interactions humaines, notamment lorsqu'ils manipulent des objects. C'est principalement dans ce dernier champ que ce sont positionnés ces travaux, avec une focalisation particulière sur les propriétés biomécaniques apparentes du membre supérieur humain.

Cette recherche a donc proposé une stratégie de contrôle interactive pour un robot, en cherchant à rendre son comportement le plus transparent possible, tout en assurant la sécurité et la stabilité de l'interaction. La méthode proposée fortement inspirée par des travaux précedents a été modifiée et complétée par d'autres techniques de manière à adapter le contrôle à une tâche précise, en se plaçant dans l'espace de cette dernière. Le contrôle a ensuite été mis en place sur un banc expérimental, pour essayer d'optimiser empiriquement le réglage de ses paramètres, et rendre le système de perturbations presque imperceptibles.

Afin de mieux comprendre la stratégie humaine lors d'une interaction avec un robot, pour la réalisation d'une tâche rythmique nécessitant un contrôle actif, une expérience a été proposée et réalisée sur une trentaine de personnes qui ont pu manipuler, pour nombre d'entre eux pour la première fois, un robot articulé. Le ressenti des participants a pu être recueilli et leurs données traitées pour identifier les caractéristiques d'un modèle très simple permettant de reproduire la dynamique de leur rejet de perturbation pendant l'interaction avec le robot. Les résultats obtenus sont cohérents avec ceux de la littérature, et viennent compléter les études comportementales des propriétés viscoélastiques cartésiennes générées par le système neuro-musculo-squelettique. 

La force d'un tel modèle réside aussi dans sa capacité d'amalgamer plusieurs phénomènes complexes émergents de sources différentes, pour engendrer des comportements globaux assimilables en grande partie à un modèle linéaire en impédance. Ce dernier facilite les analyses conduites pour étudier les performances et la robutesse des robots lorsqu'ils sont en contact physique avec des humains, comme le montre l'ébauche du dernier chapitre de ce manuscrit.
% conclusion ??

Les travaux présentés ici ont ainsi essayé de relever plusieurs des défis posés par l'estimation des propriétés biomécaniques humaines, lors de mouvements dynamiques en interaction avec un robot, en mettant en place et éprouvant différentes techniques issues de domaines variés. Plusieurs axes restent à explorer et des développements pourraient permettre d'améliorer les résultats présentés.

Comme expliqué lors de la mise en place de la méthodologie pour l'estimation des paramètres en impédance, la plus grande source d'erreurs semble provenir de l'estimation de trajectoire en position. L'amélioration de cette estimation permettrait de directement améliorer la fiabilité de la méthode. L'utilisation de réseau de neurones pourrait constituer une piste d'amélioration, cette technique ayant été implémentée avec des résultats prometteurs, toutefois limités par la base de données.

Par ailleurs, la méthode d'identification proposée suppose que les variations d'impédance ne sont pas significatives sur de faibles intervalles de temps, ce qui limite l'observation des variations cycliques des paramètres du modèle en impédance utilisé. Les méthodes non-paramétriques évoquées dans la courte revue du troisième chapitre pourraient permettre de mieux observer les variations cycliques des propriétés viscoélastiques apparentes. Pour améliorer la fiabilité de ces méthodes, les mesures des positions articulaires humaines pourraient s'avérer nécessaires, ce qui augmenterait toutefois la complexité expérimentale.

Une étude plus fine des variations d'impédance au cours des cycles d'une tâche rythmique pourrait permettre de mieux comprendre leur rôle notamment en terme de robustesse. Une méthode de contrôle en impédance robuste ou adaptative pourrait par la suite être combinée aux signaux de consignes rythmiques reposant sur des oscillateurs pour produire des mouvements cycliques bioinspirés capables d'interactions robustes avec des environnements anthropiques. Ce type de méthodes pourrait s'avérer très prometteur notamment pour les exosquelettes.